\section*{Capítulo \ref{ch:g7:stats} --- Organizar datos}

\begin{solution}{exo:g7:stats:1}
Contando cada nota:
\begin{center}
\renewcommand{\arraystretch}{1.2}
\begin{tabular}{l|ccccc|c}
nota & $9$ & $11$ & $12$ & $14$ & $15$ & total \\
\hline
frec. absoluta & $4$ & $3$ & $7$ & $3$ & $3$ & $20$ \\
\end{tabular}
\end{center}
\end{solution}

\begin{solution}{exo:g7:stats:2}
\omterm{def:g7:stats:frequency}{Frecuencias relativas}: $\frac{4}{20} = 20\,\%$, $\frac{3}{20} = 15\,\%$,
$\frac{7}{20} = 35\,\%$, $15\,\%$ y $15\,\%$. \omterm{def:g1:addition:def}{Suma}:
$20 + 15 + 35 + 15 + 15 = 100\,\%$. \checkmark
\end{solution}

\begin{solution}{exo:g7:stats:3}
Primera encuesta: $\frac{30}{50} = 0.6 = 60\,\%$. Segunda:
$\frac{45}{90} = 0.5 = 50\,\%$. La primera encuesta muestra la proporción
mayor, aunque tenga menos familias con un coche en cifras
absolutas.
\end{solution}

\begin{solution}{exo:g7:stats:4}
$\intco{5}{15}$: $8, 12, 9, 14, 11$ --- frecuencia $5$.
$\intco{15}{25}$: $23, 22, 18, 16$ --- frecuencia $4$.
$\intco{25}{35}$: $31, 27, 29, 25, 33$ --- frecuencia $5$.
$\intco{35}{45}$: $35$ --- frecuencia $1$. Total $15$. \checkmark
\end{solution}

\begin{solution}{exo:g7:stats:5}
Cuatro barras de alturas $8$, $12$, $6$ y $4$ sobre los valores $0$, $1$,
$2$ y $3$ (va bien una graduación cada $2$).
\end{solution}

\begin{solution}{exo:g7:stats:6}
Total, $30$ familias. \omterm{def:g7:stats:frequency}{Frecuencias relativas}: $\frac{8}{30}$, $\frac{12}{30}$,
$\frac{6}{30}$ y $\frac{4}{30}$. Ángulos ($\times 360^\circ$): $96^\circ$,
$144^\circ$, $72^\circ$ y $48^\circ$; \omterm{def:g1:addition:def}{suma}
$96 + 144 + 72 + 48 = 360^\circ$. \checkmark
\end{solution}

\begin{solution}{exo:g7:stats:7}
Porcentajes: verano $\frac{162}{360} = 45\,\%$; primavera
$\frac{90}{360} = 25\,\%$; otoño e invierno
$\frac{54}{360} = 15\,\%$ cada uno. De $60$ personas, el verano lo eligieron
$0.45 \times 60 = 27$.
\end{solution}

\begin{solution}{exo:g7:stats:8}
Clase A: $\frac{12}{30} = 40\,\%$. Clase B: $\frac{14}{40} = 35\,\%$.
La clase A tiene la proporción mayor de alumnos que van andando, aunque la clase B
tenga más en cifras absolutas.
\end{solution}

\begin{solution}{exo:g7:stats:9}
Con el eje empezando en $9\,700$, las tres barras tienen alturas
visibles $100$, $300$ y $400$ unidades: la última parece cuatro veces la
primera, y sugiere que las ventas se han cuadruplicado. Con un eje desde $0$, las barras
quedan casi iguales (de $9\,800$ a $10\,100$ hay una subida de un $3\,\%$):
la imagen honesta muestra unas ventas casi planas.
\end{solution}

\begin{solution}{exo:g7:stats:10}
Las \omterm{def:g7:stats:frequency}{frecuencias relativas} suman $1$: $f = 1 - 0.35 - 0.4 = 0.25$. \omterm{def:g7:stats:frequency}{Frecuencias absolutas} (total
$80$): $0.35 \times 80 = 28$; $0.4 \times 80 = 32$;
$0.25 \times 80 = 20$. Comprobación: $28 + 32 + 20 = 80$.
\end{solution}

\begin{solution}{exo:g7:stats:11}
Chicas: el $55\,\%$ de $400 = 220$; chicos: $180$. Chicas que comen en el comedor:
$0.4 \times 220 = 88$; chicos: $0.6 \times 180 = 108$. Total:
$88 + 108 = 196$ alumnos, es decir, el $\frac{196}{400} = 49\,\%$ del
colegio --- entre el $40\,\%$ y el $60\,\%$, más cerca de la tasa de las chicas
porque las chicas son más numerosas.
\end{solution}

\begin{solution}{pb:g7:stats:1}
\textbf{1.} \omterm{def:g7:stats:frequency}{Frecuencias absolutas}: $120 > 90$, gana el colegio A. Relativas:
$\frac{120}{400} = 0.30 = 30\,\%$ frente a
$\frac{90}{250} = 0.36 = 36\,\%$: gana el colegio B. El premio
debe seguir a la \omterm{def:g7:stats:frequency}{frecuencia relativa} --- el colegio B recicla una parte mayor
de lo que consume; el colegio A simplemente bebe más zumo.

\textbf{2.} Profesores: $\frac{18}{30} = 0.60 = 60\,\%$. Alumnos:
$\frac{210}{600} = 0.35 = 35\,\%$. Juntos:
$\frac{18 + 210}{630} = \frac{228}{630} \approx 0.36 =
36\,\%$. Los alumnos son veinte veces más numerosos, así que la
frecuencia conjunta se va casi por completo hacia su lado ---
el grupo grande lleva el peso grande.

\textbf{3.} Las \omterm{def:g7:stats:frequency}{frecuencias relativas} tienen que sumar $100\,\%$
(\cref{def:g7:stats:frequency}), pero
$45 + 30 + 20 + 10 = 105$: al menos un sector está mal.

\textbf{4.} \omterm{def:g7:stats:frequency}{Frecuencia absoluta} correspondiente a la relativa $0.15$:
$0.15 \times 40 = 6$. Las absolutas hasta ahora: $14 + 10 + 6 = 30$, así que
la última es $40 - 30 = 10$, con \omterm{def:g7:stats:frequency}{frecuencia relativa}
$\frac{10}{40} = 0.25$. (Comprobación: $0.35 + 0.25 + 0.15 + 0.25 =
1$.)

\textbf{5.} Ángulos: $0.40 \times 360 = 144^\circ$,
$0.35 \times 360 = 126^\circ$, $0.25 \times 360 = 90^\circ$;
y $144 + 126 + 90 = 360^\circ$: el disco cierra.

\textbf{6.} Volvió el $\frac{2.4}{10} = 0.24 = 24\,\%$ de las
papeletas.

\textbf{7.} Las guías telefónicas y los registros de vehículos recogían
sobre todo hogares acomodados, cuyo voto difería del
del país; los $2.4$ millones de respuestas eran un recuento gigantesco sacado
de la población equivocada. Como en la pregunta 1: lo que importa no
es a cuántos cuentas, sino a quiénes --- una \omterm{def:g7:stats:frequency}{frecuencia relativa} calculada sobre una
muestra distorsionada describe la muestra, no el país.

\textbf{8.} Apoyo real:
$0.7 \times 1\,000 + 0.3 \times 9\,000 = 700 + 2\,700 = 3\,400$
partidarios de $10\,000$: el $34\,\%$. La encuesta:
$0.7 \times 500 + 0.3 \times 500 = 350 + 150 = 500$ de
$1\,000$: el $50\,\%$ --- dieciséis puntos de más, porque los
ricos son la \omterm{def:g3:division:half}{mitad} de la muestra pero la décima parte del pueblo.

\textbf{9.} Respuestas: $0.4 \times 200 = 80$ descontentos y
$0.1 \times 800 = 80$ contentos, o sea, $160$ respuestas de las cuales $80$
son de descontentos: descontento medido $\frac{80}{160} = 50\,\%$. Descontento
real: $\frac{200}{1\,000} = 20\,\%$. Nadie ha mentido ---
los descontentos sencillamente responden más y la encuesta los oye
más fuerte.

\textbf{10.} Proporciones verdaderas: una décima parte ricos, o sea, $100$
entrevistas a ricos y $900$ a modestos:
$0.7 \times 100 + 0.3 \times 900 = 70 + 270 = 340$ de
$1\,000$: la encuesta reparada predice el $34\,\%$ --- la verdad de
la pregunta 8. Cincuenta mil entrevistas bien elegidas ganan a dos
millones mal elegidas.

\textbf{11.} Hospital A: leves, $\frac{90}{100} = 90\,\%$; graves,
$\frac{280}{400} = 70\,\%$; en total,
$\frac{90 + 280}{500} = \frac{370}{500} = 74\,\%$.

\textbf{12.} Hospital B: leves, $\frac{340}{400} = 85\,\%$;
graves, $\frac{65}{100} = 65\,\%$; en total,
$\frac{340 + 65}{500} = \frac{405}{500} = 81\,\%$. El hospital A
gana en los casos leves ($90 > 85$) \emph{y} en los graves
($70 > 65$).

\textbf{13.} En total, B muestra un $81\,\%$ frente al $74\,\%$ de A:
el hospital que pierde en \emph{todas} las categorías gana en el
total. (Esta inversión tiene nombre: la paradoja de Simpson.)

\textbf{14.} Las mezclas son opuestas: los pacientes de A son en su mayoría
graves ($400$ de $500$) y los de B, en su mayoría leves ($400$ de $500$).
Los casos graves se curan menos a menudo se traten donde se traten, así que la
cifra global de A queda arrastrada hacia abajo por los casos difíciles que acepta ---
la cifra global es una mezcla ponderada y los pesos difieren entre
hospitales. Un paciente grave debería elegir el hospital A: el $70\,\%$
gana al $65\,\%$ en la única fila que le concierne.

\textbf{15.} Por ejemplo: «Su comparación mezcla pacientes leves y
graves, y los dos hospitales tratan mezclas opuestas ---
separando por gravedad, el hospital A cura una proporción mayor de
\emph{los dos} tipos. Las cifras globales solo se pueden comparar cuando los
grupos que hay detrás son parecidos; si no, deciden los pesos y no la
calidad».
\end{solution}
